% Partie : théorie de la démonstration
% Chapitre : recherche de preuves
% Section : introduction

\documentclass[../../../include/open-logic-section]{subfiles}

\begin{document}

\olfileid[fr]{pt}{ps}{int}

\olsection{Introduction}

Les axiomes et les règles d'inférence de systèmes tels que
\Log{G3c} ou \Log{N1i} déterminent quels arbres de séquents ou
de formules constituent des preuves correctes. Lorsqu'un tel
arbre est donné, avec, si nécessaire, l'indication des règles
appliquées et les étiquettes des hypothèses et de leurs décharges,
ces définitions permettent de vérifier qu'il s'agit bien d'une
preuve. \emph{Trouver} une preuve d'un séquent, ou d'une formule
à partir d'hypothèses données, est un autre problème, généralement
plus difficile : celui de la \emph{recherche de preuves}.

Une procédure élémentaire consiste à énumérer les preuves.
On peut représenter les formules par des suites de symboles
sur un alphabet fini, en codant au besoin les noms et les indices.
On peut également coder les séquents et les arbres de séquents
ou de formules par de telles suites, avec des délimiteurs
pour indiquer les ramifications. Les axiomes et les règles
permettent alors de vérifier mécaniquement si une suite donnée
code une preuve correcte. Il suffit donc d'énumérer toutes
les suites finies de symboles et de ne retenir que celles
qui codent des preuves. Pour chercher une preuve du séquent
donné, on poursuit cette énumération jusqu'à en trouver une
dont il soit la conclusion. Pour une formule et des hypothèses
données, on demande de même que les hypothèses ouvertes
appartiennent à l'ensemble autorisé ; cette vérification est
mécanique lorsque l'appartenance à cet ensemble est décidable,
en particulier lorsqu'il est fini.

Une méthode plus directe reprend la démarche usuelle d'une
recherche à la main. On part du séquent final, on choisit une
formule et on applique une règle d'inférence \emph{à rebours}.
On obtient ainsi une ou plusieurs prémisses : leurs preuves,
suivies de cette inférence, donneraient une preuve du séquent
cherché. On peut procéder de manière analogue en déduction
naturelle, bien que la recherche y soit plus complexe.

Cette méthode ne garantit pourtant pas à elle seule le succès.
Un mauvais choix de règle, ou un mauvais ordre de traitement
des formules, peut conduire à une impasse. La procédure n'est
pas non plus entièrement déterminée : plusieurs formules
ou plusieurs règles peuvent être disponibles à une étape.
Un \emph{algorithme de recherche de preuves} précise tous ces
choix. Nous dirons qu'il est \emph{complet} s'il trouve une
preuve dès qu'il en existe une.

Certains systèmes se prêtent mieux que d'autres à la recherche
de preuves. Dans les calculs des séquents, l'opérateur principal
d'une formule et le membre où elle figure indiquent la règle
logique à appliquer à rebours. Ainsi, pour chercher une preuve
de $\Gamma \Sequent \Delta, !A \land !B$, où $!A \land !B$
figure dans le conséquent, on peut appliquer
\RightR{\land} à rebours et obtenir les deux prémisses
$\Gamma \Sequent \Delta, !A$ et $\Gamma \Sequent \Delta, !B$.
Dans \Log{G1c}, la contraction multiplie toutefois les choix.
On peut aussi appliquer \RightR{\Contraction} à rebours
et obtenir $\Gamma \Sequent \Delta, !A \land !B, !A \land !B$,
puis $\Gamma \Sequent \Delta, !A \land !B, !A \land !B, !A \land !B$,
et ainsi de suite.

Dans \Log{G2c}, l'application à rebours de \RightR{\land}
exige en outre de répartir les contextes : il faut choisir
$\Gamma_1$, $\Gamma_2$ tels que $\Gamma = \Gamma_1 \uplus \Gamma_2$
et $\Delta_1$, $\Delta_2$ tels que $\Delta = \Delta_1 \uplus \Delta_2$,
où l'union des multiensembles additionne les multiplicités.
On cherche ensuite des preuves de
$\Gamma_1 \Sequent \Delta_1, !A$ et
$\Gamma_2 \Sequent \Delta_2, !B$.
Une répartition mal choisie peut conduire à une impasse et
obliger à revenir en arrière pour choisir d'autres
$\Gamma_1$, $\Gamma_2$, $\Delta_1$, $\Delta_2$.
Si la formule principale est $!A \lor !B$, il faut choisir
l'une des deux règles \RightR{\lor}, donnant pour prémisse
$\Gamma \Sequent \Delta, !A$ ou $\Gamma \Sequent \Delta, !B$.
Un mauvais choix peut, là encore, imposer un retour en arrière.

Le système \Log{G3c} évite ces difficultés. Il ne comporte
pas de règles structurelles, et l'opérateur principal ainsi
que le membre où figure la formule déterminent la règle
logique à appliquer. Il se prête donc mieux à la recherche
que \Log{G1c} ou \Log{G2c}. Une recherche réussie fournit
une preuve dans \Log{G3c}, que l'on peut ensuite traduire
dans \Log{G1c}, \Log{G2c}, ou encore \Log{N1c}.
Un algorithme de recherche pour \Log{G3c}, associé à ces
traductions de preuves, fournit ainsi une méthode de recherche
pour les autres systèmes.

Comment établir la complétude d'un algorithme de recherche ?
Il faut montrer que, s'il ne trouve aucune preuve, aucune
preuve n'existe. Un algorithme mal conçu peut manquer une
preuve existante. L'énumération exhaustive évite ce défaut,
puisqu'elle examine toutes les preuves possibles, mais
elle est peu efficace. On souhaite une procédure qui réduise
le travail de recherche ; la complétude ne garantit à elle
seule ni l'efficacité ni un effort minimal.

L'idée essentielle est de tirer du déroulement d'une recherche
infructueuse un \emph{contre-modèle} du séquent initial.
En logique classique du premier ordre, $\Gamma \Sequent \Delta$
est valide si, dans toute structure, au moins une formule
de $\Gamma$ est fausse ou au moins une formule de~$\Delta$
est vraie. Le calcul \Log{G3c} est \emph{correct} :
il ne prouve que des séquents valides. On l'établit par
récurrence sur les preuves. Les axiomes sont valides
et une règle à prémisses valides a une conclusion valide.
Pour établir la complétude de la recherche, on montre que,
si elle ne trouve pas de preuve de $\Gamma \Sequent \Delta$,
on peut définir une structure dans laquelle toutes les
formules de $\Gamma$ sont vraies et toutes celles
de~$\Delta$ sont fausses. Le séquent
$\Gamma \Sequent \Delta$ est donc invalide
et, par correction du calcul, il n'a pas de preuve.
Cela établit aussi la \emph{complétude} de \Log{G3c} :
si $\Gamma \Sequent \Delta$ est valide, il possède une
preuve dans ce calcul. Toute recherche menée selon
l'algorithme complet décrit ci-dessous réussit donc sur
un séquent valide.

La procédure et le modèle de termes de ce chapitre portent
sur des séquents de formules closes, dans un langage sans
égalité logique. Les constantes servent de paramètres
et les termes d'instanciation sont clos. En présence de
variables libres ou d'égalité logique, il faut des
conventions ou des règles supplémentaires.

Une recherche infructueuse peut se poursuivre indéfiniment.
L'argument du contre-modèle établit alors l'invalidité du
séquent, même si la procédure poursuit la recherche.




\end{document}
